\section{Sensores de Temperatura LM35 y DHT11}
El sensor \textbf{LM35} es un sensor de temperatura analógico de precisión. Entrega un voltaje de salida linealmente proporcional a la temperatura en grados Celsius (con una escala de $\approx 10~\mathrm{mV}/^\circ\mathrm{C}$), facilitando su uso con conversores A/D simples \cite{TI-LM35}. A diferencia de otros sensores calibrados en unidades absolutas (Kelvin), el LM35 proporciona directamente la temperatura en Celsius sin necesidad de restar constantes de offset. Además, gracias a su calibración interna en fábrica, puede lograr una exactitud típica de $\pm 0.25\,^\circ\mathrm{C}$ a temperatura ambiente (25~$^\circ$C) sin requerir calibración externa adicional \cite{TI-LM35}. El rango de medida del LM35 abarca aproximadamente desde $-55\,^\circ\mathrm{C}$ hasta $150\,^\circ\mathrm{C}$, con una no linealidad máxima de solo $\pm 0.25\,^\circ\mathrm{C}$ y un bajo consumo de corriente (orden de 60~$\mu$A) \cite{TI-LM35}, características que lo hacen adecuado para aplicaciones de monitoreo remoto alimentadas por baterías.
\\\\
Por su parte, el \textbf{DHT11} es un sensor digital combinado de temperatura y humedad. Internamente incluye un sensor capacitivo de humedad y un termistor tipo NTC para temperatura, junto con un microcontrolador de 8 bits que realiza la conversión analógica-digital y entrega los datos ya calibrados mediante un protocolo serial de un solo hilo \cite{Aosong-DHT11}.  Este sensor proporciona directamente lecturas digitales de temperatura (en C) y humedad relativa (\%) con una calibración de fábrica que garantiza su estabilidad a largo plazo \cite{Aosong-DHT11}. El DHT11 opera típicamente en el rango de 0-50~C con una precisión aproximada de 
$\pm 2\,^\circ\mathrm{C}$, y de 20-80~\% de humedad con una precisión de $\pm 5\%$ \cite{Aosong-DHT11}. Si bien su resolución es limitada (entrega enteros: 1~C y 1~\%HR) y su frecuencia máxima de muestreo es de 1~Hz, su bajo costo y simplicidad lo han popularizado en proyectos de IoT domésticos y educativos. En comparación, existen variantes mejoradas como el DHT22 (AM2302) con mayor precisión y rango, pero el DHT11 sigue siendo adecuado cuando se busca una solución económica para medir condiciones ambientales básicas.
\section{Protocolo MQTT y Arquitectura Cliente-Servidor}
MQTT es un protocolo de comunicación orientado a mensajería publicado/suscrita que opera sobre TCP/IP, diseñado específicamente para dispositivos con recursos limitados en entornos IoT. Fue originalmente desarrollado por IBM en 1999 y estandarizado por OASIS en 2014, convirtiéndose desde entonces en un pilar para aplicaciones IoT por su eficiencia y simplicidad
mqtt.org \cite{awsMQTT}. A diferencia de un modelo HTTP tradicional (cliente-servidor estricto), MQTT implementa una arquitectura de \textbf{publicación/suscripción} mediada por un servidor central denominado \textit{broker} o agente de mensajes. En este esquema, los dispositivos actúan como clientes MQTT que se conectan al broker y pueden desempeñar dos roles lógicos: \textit{editores} (publishers), que envían mensajes sobre ciertos \textit{topics} (temas); y \textit{suscriptores} (subscribers), que reciben mensajes de los topics a los que se hayan suscrito
\cite{awsMQTT}. El broker MQTT central es el encargado de recibir todos los mensajes publicados por los clientes editores y distribuirlos a los clientes suscritos a cada topic correspondiente
\cite{awsMQTT}. Este desacoplamiento espacial y temporal entre remitentes y receptores aporta gran flexibilidad: los clientes no necesitan conocerse entre sí ni estar activos simultáneamente, solo deben mantener una conexión con el broker. La comunicación mediante MQTT se realiza de forma eficiente en cuanto a cabecera y tamaño de los mensajes. Un mensaje MQTT típico consta de un tópico (cadena de texto jerárquica que identifica el tema) y un payload de datos binarios de longitud reducida. El protocolo define tres niveles de \textbf{calidad de servicio} (QoS 0, 1 y 2) para garantizar la entrega de mensajes según la criticidad de la aplicación
\cite{awsMQTT}. Por ejemplo, en QoS0 el mensaje se entrega “al menos una vez” sin acuse de recibo, mientras que en QoS2 se asegura entrega “exactamente una vez” mediante confirmaciones y reenvíos, evitando duplicados. Asimismo, MQTT soporta sesiones persistentes que facilitan la reconexión de clientes en redes móviles inestables, y permite el uso de mecanismos de seguridad (TLS para cifrado, autenticación mediante usuarios/OAuth, etc.)
\cite{mqttOrg}. Todas estas características han contribuido a que MQTT sea ampliamente adoptado en plataformas IoT comerciales y proyectos de investigación, siendo compatible con servicios en la nube (AWS IoT, Azure IoT, IBM Cloud, etc.) y una gran variedad de dispositivos embebidos.
\section{Comunicación Serial UART entre Microcontroladores}
En los experimentos desarrollados, la conexión entre el Arduino Mega y el ESP32 se realiza a través de una interfaz serial UART. La \textbf{UART} (\textit{Universal Asynchronous Receiver-Transmitter}) es uno de los métodos más sencillos y difundidos para la comunicación directa entre microcontroladores
\cite{Analog-UART}. Se trata de un protocolo de comunicación serial \textit{asíncrono}, lo que significa que no utiliza una señal de reloj compartida entre emisor y receptor; en su lugar, ambos dispositivos deben configurarse a la misma velocidad de transmisión (baudrate) para lograr sincronización en la recepción de bits
\cite{Analog-UART}. Una UART transmite los datos bit a bit (serialmente) a través de una línea de transmisión (TX) mientras el receptor escucha en su línea RX opuesta. De este modo, con solo un par de líneas (TX y RX cruzadas entre dos dispositivos, más una referencia común de tierra) es posible implementar un canal de comunicación full-duplex simple
\cite{Analog-UART}.
\\\\
La simplicidad de la UART la hace muy atractiva para interconectar componentes en un sistema embebido heterogéneo. Muchos microcontroladores, incluyendo el ATmega2560 (Arduino Mega) y el ESP32, incorporan módulos UART hardware integrados, lo que permite establecer puentes de comunicación serial con apenas configurar el baudrate y parámetros de trama (p.ej., 8 bits de datos, sin paridad, 1 bit de stop: “8N1”). En una arquitectura IoT como la descrita, la UART sirve para que el Arduino Mega envíe los datos de temperatura leídos por los sensores al ESP32. El ESP32, a su vez, actúa como nodo de comunicaciones inalámbricas, publicando esos datos mediante WiFi al broker MQTT. Esta división de tareas aprovecha las fortalezas de cada microcontrolador (el Mega en adquisición de datos mediante sus múltiples pines analógicos/digitales, y el ESP32 en conectividad de red), logrando un sistema modular. La comunicación UART entre ambos se configura típicamente a un baudrate adecuado (por ejemplo 9600 o 115200 bps) para asegurar una transferencia de datos suficientemente rápida respecto al periodo de muestreo de los sensores, manteniendo errores de sincronización despreciables. Cabe resaltar que, por ser un enlace asíncrono de nivel TTL, es importante compartir la referencia de tierra entre los dispositivos y, si operan a diferentes voltajes lógicos (5V en Arduino Mega vs 3.3V en ESP32), utilizar un adaptador de nivel lógico para evitar daños en los puertos seriales. En resumen, el uso de UART proporciona una forma fiable y de baja complejidad para la comunicación punto a punto entre microcontroladores en prototipos IoT, facilitando el intercambio de datos sensor y control en tiempo real.